\section{Case Study: Fuel Crisis and Tax Rebate}

\begin{frame}{Fuel Crisis and Fuel Tax Cut}
    \showmytikz[1]{../shared/tikz/dia-tikz-petrolprices.tex}
    \mycomment{The war in Iran lead to an increase of the prices for petrol and diesel in Germany. There was considerable delay due to tankers still on their way. 
    
    The German Federal Government reduced the taxes per liter petrol by 17 cent by May 1st, but prices went down less. Given inelastic demand and elastic supply, it was expected the that the price would drop almost in the full amount of the tax rebate. This did not happen.
    
    A similar policy was pursued in Canada with a fuel excise tax cut of 10ct from April 20th, also to no avail.---Note that the difference between the end price is nearly completely due to taxation; both government impose a number of per unit and ad valorem taxes, with much higher taxation in Germany.
    }
\end{frame}

\begin{frame}{Fuel Crisis and Fuel Tax Cut}
    \mycomment{
        If this happens in the more typical situation where demand is highly inelastic while supply is pretty elastic, the restriction can have a significant welfare effect.

        An example would be the forced reduction of world wide available crude oil in 2026.
    }
    \showmytikz[1,2]{../shared/tikz/tikz-quantity-oil}
\end{frame}

\begin{frame}{Fuel Crisis and Fuel Tax Cut}
    \mycomment{
        Consumers lose out and producers win - that is, those producers who are in the market still. And all in all, welfare is reduced relative to status quo ante.

        This is because the new situation might move quantities away from the competitive equilibrium towards the Cournot point, where producer surplus is maximum.

        In this situation the temptation to tax those profits is indeed considerable. That would shift the supply curve up, with almost no influence on price. It is, however, difficult to discern where exactly in the supply chain the profits have been realised.
    }
    \showmytikz[2,3]{../shared/tikz/tikz-quantity-oil}
\end{frame}

\begin{frame}{Fuel Crisis and Fuel Tax Cut}
    \mycomment{
        Imagine the state dropping a unit tax on the now restricted good. The effect of the tax drop on the price ($p_2-p_3$) is (much) lower than the tax rebate $\Delta t$.
    }
    \showmytikz[2,3,4]{../shared/tikz/tikz-quantity-oil}
\end{frame}
\begin{frame}{Fuel Crisis and Fuel Tax Cut}
    \mycomment{
        The change to consumer surplus will be small, but the producer surplus will get larger.

    }
    \showmytikz[3,4,5]{../shared/tikz/tikz-quantity-oil}
\end{frame}

\begin{frame}{Fuel Crisis and Fuel Tax Cut}
    \mycomment{
        The loss of government revenue ($\Delta T \cdot x_2$) is larger than the reduction of the dead weight loss. Both consumers and producers gain; who exactly how much will depend on the ratio of elasticities.

    }
    \showmytikz[1,4,5,6]{../shared/tikz/tikz-quantity-oil}
\end{frame}


